% Generated draft: requires mathematical and source review.
\documentclass[11pt]{article}
\usepackage[margin=25mm]{geometry}
\usepackage{fontspec}
\setmainfont{Latin Modern Roman}
\usepackage{amsmath,amssymb,mathtools,unicode-math}
\setmathfont{Latin Modern Math}
\usepackage{enumitem,needspace,xurl,hyperref}
\setlength{\parindent}{0pt}
\setlength{\parskip}{5pt}
\setlength{\emergencystretch}{3em}
\newcommand{\N}{\mathbb{N}}
\newcommand{\Z}{\mathbb{Z}}
\newcommand{\Q}{\mathbb{Q}}
\newcommand{\R}{\mathbb{R}}
\newcommand{\C}{\mathbb{C}}
\newcommand{\F}{\mathbb{F}}
\newcommand{\A}{\mathbb{A}}
\newcommand{\PP}{\mathbb{P}}
\newcommand{\E}{\mathbb{E}}
\newcommand{\eps}{\varepsilon}
\newcommand{\dd}{\,\mathrm{d}}
\begin{document}
\section*{UnsolvedMath ID: 30000507}
\textbf{Source problem code:} catalogue entry
\subsection{Definitions and mathematical statement}
The source asks which types of algebraic curves occur as Teichm\"uller curves. In the usual unpointed genus-$g\geq2$ setting, a nonzero holomorphic quadratic differential on a marked Riemann surface determines a Teichm\"uller disk $\mathbb H\to\mathcal T_g$, isometric for the relevant Teichm\"uller metric. A Teichm\"uller curve is an algebraic curve arising from its projection to the moduli space $\mathcal M_g$, with the appropriate normalization/orbifold convention. Translation-surface examples use squares of abelian differentials; general quadratic differentials need not be squares.

For a flat surface, the associated Veech group is the derivative group of its affine automorphisms, viewed in the appropriate $\operatorname{PSL}_2(\R)$ convention. The lattice condition, finite hyperbolic covolume, provides the usual algebraic Teichm\"uller-curve criterion. The resulting normalized curve is described by an orbifold quotient $\mathbb H/\Gamma$ in that framework. This recognition criterion for a specified flat surface does not classify the abstract algebraic curves that can occur.

A precise classification request must distinguish: abstract curves versus their maps/images in moduli space, orbifold curves versus coarse curves, punctured curves versus smooth projective compactifications, and actual normalizations versus arbitrary finite covers. It must also fix whether genus/marked points and strata are prescribed or allowed to vary. The source does not choose these meanings of ``types'', and no restriction to translation surfaces, a fixed genus or a named family is silently imposed.

The research target is to describe the possible algebraic-curve types in the intended convention, with existence and exclusion criteria or a full classification as appropriate. The supplied literature triage checked 21 August 2026 reports Veech-group criteria and many families but says no complete general classification answers the broad question. This is a dated supplied assessment, not an independent current literature determination.
\subsection{Short English statement}
Which algebraic curves can be Teichm\"uller curves? Lattice Veech groups recognize many examples, but the question still needs a convention about abstract curves, orbifolds, compactifications, covers and ambient genus. The supplied assessment reports no complete general classification.
\subsection{Sources}
\begin{itemize}
\item Local UnsolvedMath registry, permanent ID 30000507, exact source code catalogue entry; \emph{Classification of Algebraic Teichm\"uller Curves}. Proposed year 2006, status open; exact broad classification question, with supplied verification metadata reporting normalization/checking.
\item Supplied report citation: \emph{Teichm\"uller Space (Classical and Quantum) (2007)}, Oberwolfach Reports, DOI 10.4171/owr/2006/26, pp. 1537--1614, \url{https://doi.org/10.4171/owr/2006/26}. Title year differs from DOI/proposed year.
\item Supplied literature triage checked 2026-08-21: Veech, \emph{Teichm\"uller curves in moduli space}, \url{https://doi.org/10.1007/BF01388890}; ``original Oberwolfach report'', \url{https://doi.org/10.4171/owr/2006/26}. This dated OWR-derived curation requests independent verification of time-sensitive status and cited claims.
\end{itemize}
\end{document}
